\chapter{Discussion}
\label{ch:discussion}

\section{Interpreting the Detection Performance}

The Stage 1 evaluation results present a nuanced picture of what multi-agent LLM-based intrusion detection can and cannot achieve at the current frontier of model capability. Across thirteen of the fourteen attack types in the CICIDS2018 benchmark, AMATAS achieves a mean F1 score of 92.9 per cent and a mean recall of 90.2 per cent, with a false positive rate of zero per cent in all thirteen cases. These figures position AMATAS in a distinctive location within the landscape of existing NIDS research --- not the highest raw accuracy ever reported, but occupying a performance envelope that no prior system combines with systematic, analyst-readable reasoning.

Direct comparison with related work is instructive. \textcite{Houssel2024}, in the closest prior study, evaluated a standalone LLM-based NIDS and reported detection rates of approximately fifty per cent on benchmark flows, with significant false positive inflation. AMATAS achieves roughly double that recall, a gain attributable to the shift from single-agent to multi-agent architecture, the adversarial pressure of the Devil's Advocate, and the temporal context supplied by the cross-flow correlation agent. At the other extreme, \textcite{GutierrezGaleano2025} demonstrated F1 scores of 99.84 per cent using a fine-tuned T5 language model, but their approach produces no reasoning whatsoever --- the fine-tuned model maps feature vectors to labels through learned weights that are as opaque as any traditional classifier. \textcite{Li2024} present IDS-Agent, a ReAct-style single-agent system reporting an F1 of 0.75 on the CIC-IoT dataset; AMATAS achieves F1=0.88 on a more comprehensive benchmark whilst additionally generating the reasoning chains that IDS-Agent does not systematically produce.

The appropriate analytical frame for these comparisons is not whether AMATAS outperforms fine-tuned discriminative models on raw accuracy --- it does not, and was not designed to. The relevant question is whether AMATAS achieves accuracy that is operationally useful whilst simultaneously providing transparent, human-readable reasoning that fine-tuned models structurally cannot. On that framing, the results are strongly affirmative. The comparison with \textcite{Kalafatidis2025}, who independently propose a two-tier LLM-enhanced IDS architecture, further situates AMATAS within a converging research consensus --- though AMATAS extends that pattern with the multi-agent debate structure and systematic explainability artefacts that prior two-tier proposals do not include.

Within the Stage 1 results, the performance profile across attack types reveals important structural patterns. Brute-force and volumetric DoS categories --- FTP-BruteForce (F1=100\%), SSH-Bruteforce (F1=99\%), DoS-Slowloris (F1=100\%), DoS-SlowHTTPTest (F1=100\%) --- achieve the highest scores, reflecting the distinctive NetFlow signatures these attacks produce. DDoS-HOIC (F1=72\%, recall=58\%) represents the hardest detectable category and is discussed at length in Section~\ref{sec:hoic}. Infiltration (F1=0\%, recall=0\%) is treated separately as a fundamental architectural limitation rather than a conventional performance shortfall.

\section{The Explainability Advantage}

The central thesis of this work is not that LLMs detect intrusions more accurately than traditional classifiers --- on raw accuracy alone, this claim would be false for well-tuned discriminative models. The thesis is that LLMs can detect intrusions with operationally viable accuracy whilst producing explanations of a quality and depth that traditional classifiers structurally cannot. Every AMATAS detection produces six distinct analytical outputs.

Consider flow 783 from the FTP-BruteForce batch as a representative example. The Protocol Agent classifies the flow as BENIGN at confidence 0.90, finding that the port assignment and protocol combination are individually consistent with legitimate FTP traffic. The Statistical Agent escalates to SUSPICIOUS at 0.75, flagging anomalous byte-per-packet ratios and inter-arrival timing. The Behavioural Agent returns SUSPICIOUS at 0.80, identifying a pattern consistent with credential stuffing but insufficient alone to confirm it. The Temporal Agent delivers MALICIOUS at confidence 0.95, having correlated the flow with 49 other connections from the same source IP in the analysis window --- 49 connections to port 21 from a single source within a compressed time window is the decisive evidence that no single-flow analysis could provide. The Devil's Advocate acknowledges the temporal cluster but argues that the source IP could plausibly belong to a legitimate FTP client with an aggressive retry policy. The Orchestrator synthesises all five analyses into a SUSPICIOUS verdict at confidence 0.65 with consensus score 0.75, maps the finding to MITRE ATT\&CK technique T1110 (Brute Force), and produces a narrative explicitly reconciling the Protocol Agent's benign reading with the Temporal Agent's cluster evidence.

This output is qualitatively different from what SHAP or LIME explanations provide for traditional classifiers. SHAP values tell an analyst which features influenced a model's output; they do not articulate what those features mean in the context of the specific flow, how they relate to known attack patterns, whether a reasonable benign interpretation exists, or how competing evidence was weighed. AMATAS produces all of these, in natural language accessible to analysts without machine learning expertise. This aligns with the chain-of-thought reasoning framework established by \textcite{Wei2022}, who demonstrate that reasoning traces generated during inference are independently verifiable in a way that post-hoc attributions are not, and with \textcite{Hesford2024}, who argue that the most productive analyst-NIDS relationship is one where the system surfaces structured evidence rather than opaque verdicts. The practical implication is that AMATAS supports analyst decision-making rather than displacing it: an analyst reviewing a verdict can assess the reasoning, identify which agent analysis they find unconvincing, and apply domain expertise to confirm or override the system's conclusion.

\section{The False Positive Problem Resolved}

The zero per cent false positive rate across thirteen of fourteen attack types is perhaps the most practically significant result of this evaluation. Alert fatigue --- the systematic desensitisation of security analysts to alerts generated by over-sensitive detection systems --- is one of the most well-documented operational problems in network security. \textcite{Alahmadi2020} report that analysts in high-volume environments routinely ignore substantial fractions of generated alerts, with the consequence that genuine threats embedded within alert storms go uninvestigated. A NIDS that maintains near-zero false positive rates is not merely more accurate; it is operationally usable in a way that a high-FPR system is not, regardless of the latter's recall advantage. At the CICIDS2018 dataset scale of twenty million flows, even a one per cent false positive rate would generate two hundred thousand spurious alerts --- a volume that would render any system operationally useless.

The mechanism responsible for AMATAS's false positive suppression is the Devil's Advocate agent operating at a weight of thirty per cent in the Orchestrator's consensus calculation. This weight was calibrated through empirical ablation: at fifty per cent (Phase 3e), the agent became excessively aggressive, depressing recall below acceptable thresholds; at thirty per cent, it applies sufficient adversarial pressure to prevent the specialist agents from over-classifying ambiguous flows as malicious, without overwhelming genuinely malicious signals. This finding provides the first demonstration of the adversarial agent mechanism, described conceptually by \textcite{Yin2024} and the multi-agent debate literature \autocite{Du2023, Chan2023}, applied specifically to network intrusion detection with quantified performance consequences across a fourteen-category evaluation.

\section{Limitations and Failure Modes}
\label{sec:limitations}

\subsection{The Infiltration Failure}

The Infiltration category produces a complete detection failure: zero true positives from fifty attack flows. This is not a performance shortfall addressable through prompt engineering or model selection; it represents a fundamental limitation of flow-level analysis. Infiltration attacks in the CICIDS2018 dataset are implemented via a simulated remote code execution chain that, at the individual flow level, generates traffic indistinguishable from normal HTTPS or SMB communication. There is no anomalous packet rate, no unusual port combination, no volume spike --- the attacking flows conform in every measurable dimension to the statistical profile of the benign traffic in which they are embedded.

The solution this failure motivates is temporal clustering, the v3 architectural extension described in Section~\ref{sec:future}. The hypothesis is that providing the agent pipeline with ordered sequences of flows from the same source IP, rather than individual flows, would reveal the sequential access progression characterising infiltration attacks. Individual flows from an infiltrating host look benign; the ordered sequence of those flows --- accessing a progression of internal services over a compressed time window --- may not.

\subsection{DDoS-HOIC and Feature Overlap}
\label{sec:hoic}

DDoS-HOIC achieves recall of 58 per cent and F1 of 72 per cent --- the lowest among successfully detected categories. The High Orbit Ion Cannon tool generates HTTP floods using randomised headers, user agents, and URL paths specifically to evade signature-based detection. The resulting flows share characteristics with high-volume legitimate web traffic: they use standard ports, conform to expected HTTP packet size distributions, and do not exhibit the periodicity that characterises simpler volumetric attacks. The Statistical Agent correctly identifies anomalous flow rates in 58 per cent of cases, but in the remaining 42 per cent the HOIC randomisation pushes the flow's statistical profile into the range occupied by legitimate web servers under genuine load.

\subsection{MCP Tool Utility and Anonymised Datasets}

The MCP comparison experiment provides a controlled test of the hypothesis that external threat intelligence tools provide no benefit on anonymised IP address spaces. The results confirm it precisely. Configuration C (engineered prompt with MITRE ATT\&CK MCP tool) achieves F1=58.3 per cent versus F1=58.0 per cent for Configuration B (engineered prompt alone) --- a difference of 0.3 percentage points at additional cost. Tools querying AbuseIPDB, AlienVault OTX, and geolocation services return no actionable information on RFC 1918 private IP addresses that constitute the entire CICIDS2018 address space. This validates the AMATAS design choice of embedding domain knowledge in agent prompts rather than in external tool calls, and confirms the theoretical expectation articulated for tool-augmented LLM systems \autocite{Schick2023, Yao2023, Qin2024}: when tools return null or irrelevant results, tool-augmented architectures are strictly inferior to prompt-only architectures, incurring additional latency and cost without informational gain.

The 25.2 percentage-point F1 gap between the best single-agent configuration (Config A, F1=62.8\%) and AMATAS (F1=88\%) quantifies the specific value added by multi-agent debate, specialist role assignment, and adversarial argumentation. This gap is not an artefact of model selection --- Config A uses GPT-4o-mini whilst AMATAS uses GPT-4o --- but a structural consequence of the architectural differences: four specialist perspectives, adversarial counter-argumentation, and Orchestrator synthesis collectively overcome the reasoning ceiling observable in any single-agent configuration.

\subsection{Model Selection Constraints and Generalisability}

The exclusive reliance on GPT-4o means that reported performance figures may not generalise to other frontier LLMs. GPT-4o-mini exhibited a systematic false positive bias --- classifying every flow as malicious regardless of content --- which rendered it unusable for NIDS despite its cost advantage. Claude Sonnet 3.5 was rejected due to rate limits incompatible with 1000-flow batch processing and an estimated cost of approximately \$108 per batch. The multi-agent architecture is model-agnostic in principle, but the specific calibration of agent prompts, confidence thresholds, and Devil's Advocate weights was developed against GPT-4o behaviour. Alternative models may exhibit different reasoning patterns, hallucination profiles, and confidence tendencies requiring re-calibration.

The CICIDS2018 dataset, whilst the most comprehensive publicly available benchmark \autocite{Sharafaldin2018, Leevy2020}, is a simulation of network traffic generated in a controlled laboratory environment. The FLOW\_DIRECTION feature was unavailable in the CICIDS2018 NetFlow v3 encoding and was excluded from the Stage 1 feature set, which may limit discriminative power on directionally distinctive attacks. \textcite{Sarhan2021} note several limitations of the NetFlow v3 representation, and \textcite{FengSakurai2025} contextualise dataset representativeness limitations within the broader challenge of constructing reliable NIDS evaluation benchmarks.

\section{Security Considerations for MCP-Enhanced Detection Systems}
\label{sec:mcp-security}

The MCP comparison experiments demonstrate that external tool access provides negligible benefit on anonymised datasets, but a production deployment of MCP-enhanced intrusion detection against live network traffic would face a qualitatively different threat model. The security properties of the MCP architecture itself --- the protocol through which an LLM agent discovers, invokes, and trusts external tools --- become critical concerns when the system operates on real traffic with genuine IP addresses, domain names, and payload content.

The most immediate risk is tool injection, where a compromised or malicious MCP server returns fabricated results designed to influence the LLM agent's classification. An attacker who controls an upstream threat intelligence feed could supply false reputation data --- reporting a known malicious IP address as benign, or conversely, poisoning the reputation of legitimate infrastructure to trigger false positives at scale. The LLM agent, lacking independent means to verify the provenance or integrity of tool responses, would incorporate the fabricated data into its reasoning as though it were authoritative. This vulnerability is structural rather than incidental: the MCP specification delegates trust to the tool server, and no mechanism within the current protocol provides for cryptographic verification of response integrity or cross-validation against independent sources. \textcite{Hasan2025} provide a comprehensive analysis of security vulnerabilities in MCP-based architectures, identifying tool injection, permission escalation, and data exfiltration as the principal attack surfaces, and proposing a defence-in-depth framework that includes response signing, tool sandboxing, and anomaly detection on tool output distributions.

Permission escalation presents a related concern. An MCP server designed to provide read-only access to threat intelligence databases could, through crafted tool descriptions or parameter manipulation, induce the LLM agent to execute actions beyond its intended scope --- writing to configuration files, modifying firewall rules, or exfiltrating sensitive flow data to external endpoints. The risk is amplified in intrusion detection contexts where the system necessarily processes sensitive network telemetry, and where an escalation from analysis to action could enable an attacker to manipulate the very defences the system is meant to provide.

Data exfiltration through MCP tool calls represents a third category of risk. Each tool invocation transmits a structured request to an external server, and in a NIDS context these requests may contain IP addresses, port numbers, flow statistics, and other telemetry that an adversary could aggregate to map the defended network's topology. Even if individual tool calls reveal limited information, the aggregate of thousands of calls across an analysis session could provide substantial intelligence to a patient adversary controlling or monitoring the tool server.

These considerations do not invalidate the MCP architecture for production deployment, but they establish that any future integration of external tools into the AMATAS pipeline must be accompanied by a security layer that the current MCP specification does not provide. The AMATAS design decision to embed domain knowledge in agent prompts rather than in external tool calls is, in retrospect, not merely a consequence of the anonymised dataset's incompatibility with threat intelligence lookups, but a security-conservative architectural choice that eliminates the tool-mediated attack surface entirely for the current deployment context.

\section{Threats to Validity}

With respect to internal validity, each attack-type batch is constructed from a single random seed, meaning that performance figures reflect one specific sample of flows per category. A different seed might produce a batch with a different distribution of edge cases, potentially altering recall figures for categories with moderate sample-level variance. Future work should conduct multi-seed evaluations to establish confidence intervals.

With respect to construct validity, AMATAS produces three-class verdicts --- BENIGN, SUSPICIOUS, and MALICIOUS --- but the evaluation framework collapses SUSPICIOUS and MALICIOUS into a single detected category for binary metric calculation. This collapse is justified operationally --- a SUSPICIOUS verdict triggers investigation just as a MALICIOUS verdict does --- but it means that reported recall figures include flows classified at relatively low confidence, and a stricter MALICIOUS-only evaluation would produce lower recall figures.

The economic projection assumes that the Tier 1 Random Forest's one-hundred-per-cent recall on the development partition generalises to evaluation batches, and that cost per LLM call remains constant across flow types. In practice, the Temporal Agent's cost varies with the number of co-IP flows available for context, and flows with more complex feature patterns generate longer reasoning chains with higher token counts.
